\section{频率模型与严重度模型}

\subsection{建模对象}

项目把建模对象拆成两部分：
\begin{itemize}[leftmargin=2.2em]
  \item \textbf{频率模型}：拟合事件数量，包括总延误航班数、取消数、备降数，以及 internal/system/external 三类事件数；
  \item \textbf{严重度模型}：拟合单次事件或单位事件的平均损失，包括总延误航班平均分钟、internal/system/external 平均事件分钟。
\end{itemize}

这样的拆分符合典型 operational risk 建模思路。总体损失不是被直接拟合，而是被分解为
\begin{equation}
L = \sum_{k=1}^{N} X_k,
\end{equation}
其中 $N$ 是频率随机变量，$X_k$ 是严重度随机变量。项目虽然在最后还加入取消和备降 penalty，但整体思想仍然是标准的 \textbf{frequency--severity decomposition}。

\subsection{候选频率分布：Poisson 与 Negative Binomial}

\paragraph{Poisson 分布}

若某类事件数 $N$ 服从 Poisson 分布，则
\begin{equation}
\mathbb{P}(N=n) = \frac{e^{-\lambda}\lambda^n}{n!}, \qquad n=0,1,2,\dots
\end{equation}
其均值与方差相等：
\begin{equation}
\mathbb{E}[N]=\lambda, \qquad \mathrm{Var}(N)=\lambda.
\end{equation}

Poisson 模型适合描述随机、独立、方差与均值同量级的计数过程。

\paragraph{Negative Binomial 分布}

若计数存在明显过度离散，即方差显著大于均值，则更适合用 Negative Binomial：
\begin{equation}
\mathbb{P}(N=n) = \binom{n+r-1}{n} p^r (1-p)^n, \qquad n=0,1,2,\dots
\end{equation}
其中
\begin{equation}
\mathbb{E}[N] = r\frac{1-p}{p}, \qquad
\mathrm{Var}(N) = r\frac{1-p}{p^2}.
\end{equation}

脚本并没有通过通用数值优化直接做 MLE，而是利用样本均值 $\bar{n}$ 与样本方差 $s^2$ 做矩估式参数恢复：
\begin{align}
r &= \frac{\bar{n}^2}{s^2-\bar{n}}, \\
p &= \frac{r}{r+\bar{n}},
\end{align}
当 $s^2 \le \bar{n}$ 时，脚本把分散参数设成一个极大值，退化到近似 Poisson 的情况。

\paragraph{对数似然与 AIC}

对任意候选模型 $M$，项目计算样本对数似然
\begin{equation}
\ell(M) = \sum_{i=1}^{m}\log f_M(x_i),
\end{equation}
并使用
\begin{equation}
\mathrm{AIC}(M) = 2k - 2\ell(M)
\end{equation}
进行比较，其中 $k$ 为参数个数。AIC 越小，模型在拟合优度与复杂度之间的综合表现越好。

\subsection{频率模型结果}

项目结果表明，\textbf{所有频率变量最终都选择了 Negative Binomial}。这与数据的过度离散特征一致，例如：
\begin{itemize}[leftmargin=2.2em]
  \item 总延误航班数均值为 241.21，但方差达到 91{,}192.93；
  \item 取消数均值为 32.04，但方差达到 4{,}403.52；
  \item 备降数均值为 3.29，但方差达到 29.83。
\end{itemize}

\begin{table}[H]
\centering
\caption{选中的频率模型}
\begin{tabular}{lccrr}
\toprule
变量 & 选中分布 & AIC & normal 状态均值 & disrupted 状态均值 \\
\midrule
total\_delay\_count & Negative Binomial & 7{,}290.82 & 207.13 & 340.87 \\
cancellation\_count & Negative Binomial & 4{,}463.27 & 14.15 & 84.36 \\
diversion\_count & Negative Binomial & 2{,}531.80 & 2.60 & 5.30 \\
internal\_ops\_count & Negative Binomial & 6{,}766.28 & 132.65 & 204.68 \\
system\_ops\_count & Negative Binomial & 6{,}047.51 & 68.42 & 124.20 \\
external\_ops\_count & Negative Binomial & 3{,}502.99 & 6.06 & 12.00 \\
\bottomrule
\end{tabular}
\end{table}

状态均值的差异也很重要。disrupted 状态下，总延误事件数、取消数和 system/internal 风险块事件数都显著高于 normal 状态，这为后续 Markov 状态分析提供了统计支持。

\begin{figure}[H]
  \centering
  \includegraphics[width=0.88\textwidth]{chart_8_frequency_aic_comparison_en.png}
  \caption{频率模型 AIC 比较}
  \label{fig:freq-aic}
\end{figure}

\subsection{候选严重度分布：Lognormal 与 Weibull}

\paragraph{Lognormal 分布}

若 $\ln X \sim \mathcal{N}(\mu,\sigma^2)$，则 $X$ 服从 Lognormal 分布，其密度为
\begin{equation}
f(x) = \frac{1}{x\sigma\sqrt{2\pi}}
\exp\left(
  -\frac{(\ln x - \mu)^2}{2\sigma^2}
\right), \qquad x>0.
\end{equation}

Lognormal 适合描述右偏、长尾、严格为正的损失变量，是 operational risk 中常见的严重度候选分布。

\paragraph{Weibull 分布}

Weibull 分布的密度为
\begin{equation}
f(x) = \frac{k}{\lambda}
\left(\frac{x}{\lambda}\right)^{k-1}
\exp\left[-\left(\frac{x}{\lambda}\right)^k\right], \qquad x>0,
\end{equation}
其中 $k$ 为形状参数，$\lambda$ 为尺度参数。Weibull 在拟合寿命型、单峰、正值变量时也很常见。

\paragraph{拟合与选择}

脚本使用 \texttt{scipy.stats} 的 \texttt{fit(..., floc=0)} 对 Lognormal 与 Weibull 进行拟合，并同样使用 AIC 比较。由于严重度变量只在正值样本上有意义，因此项目先过滤掉零值，再对正样本做拟合：
\begin{equation}
\mathcal{D}^{+} = \{x_i : x_i > 0\}.
\end{equation}

\subsection{严重度模型结果}

严重度模型的选择不是单一模式，而是变量相关的：
\begin{table}[H]
\centering
\caption{选中的严重度模型}
\begin{tabular}{lccrr}
\toprule
变量 & 选中分布 & AIC & normal 状态均值 & disrupted 状态均值 \\
\midrule
total\_avg\_delay\_minutes & Weibull & 5{,}146.84 & 67.47 & 76.76 \\
internal\_avg\_delay\_minutes & Lognormal & 5{,}362.19 & 78.28 & 83.24 \\
system\_avg\_delay\_minutes & Lognormal & 4{,}413.72 & 47.83 & 65.53 \\
external\_avg\_delay\_minutes & Lognormal & 4{,}864.21 & 103.71 & 111.83 \\
\bottomrule
\end{tabular}
\end{table}

可以看到：
\begin{itemize}[leftmargin=2.2em]
  \item 总体延误航班平均分钟更适合用 Weibull；
  \item 三个风险块的单位事件严重度更适合用 Lognormal；
  \item external 风险块虽然频率最低，但单位事件平均分钟却是三类中最高的，这说明 external shock 更像低频高损的尾部来源。
\end{itemize}

\begin{figure}[H]
  \centering
  \includegraphics[width=0.88\textwidth]{chart_9_severity_aic_comparison_en.png}
  \caption{严重度模型 AIC 比较}
  \label{fig:sev-aic}
\end{figure}

\subsection{本阶段的解释}

到这一步，项目已经完成了 frequency--severity 分解中最核心的分布选择工作。结果可以概括为：
\begin{enumerate}[leftmargin=2.2em]
  \item \textbf{频率侧存在明显过度离散}，因此 Negative Binomial 全面胜出；
  \item \textbf{严重度侧没有统一赢家}，说明不同风险块的损失生成机制并不完全相同；
  \item \textbf{disrupted 状态同时抬高频率与严重度均值}，为后续引入状态依赖与 Markov 机制提供了经验基础。
\end{enumerate}
